\color{unimain}{\subsection{Блокировка при неисправности цепей напряжения}}
\color{black}

\begin{enumerate}

\item Общие сведения
\begin{enumerate}

\item
Для алгоритмов РЗА, которые при неисправностях в цепях основного ТН могут излишне или ложно сработать и повлиять на режим работы первичной сети, используется специальная блокировка при неисправностях в цепях напряжения (БНН), которая вводится программным ключом \textbf{«Контр ЦН»}. 

\end{enumerate}

\item БНН по критерию пофазного сравнения напряжений (БНН-ф)
\begin{enumerate}

\item
В качестве основного алгоритма БНН используется критерий пофазного сравнения напряжений основной вторичной обмотки, собранной по схеме «звезда» (ВО-1), и дополнительной вторичной обмоткой, собранной по схеме «разомкнутый треугольник» (ВО-2). Между замерами цепей «звезды» и приведенных к ним замерами «разомкнутого треугольника» производится расчёт векторных невязок по следующим формулам:

\begin{equation}\label{eq:BNN_a}
U_{\textsl{БНН ф.A}} = \left| \dot{U}_{A\textsl{(ВО-1)}} - \frac{\dot{U}_{a\textsl{(ВО-2)}}}{\sqrt{3}} \right|
\end{equation}

\begin{equation}\label{eq:BNN_b}
U_{\textsl{БНН ф.В}} = \left| \dot{U}_{B\textsl{(ВО-1)}} - \frac{\dot{U}_{b\textsl{(ВО-2)}}}{\sqrt{3}} \right|
\end{equation}

\begin{equation}\label{eq:BNN_c}
U_{\textsl{БНН ф.С}} = \left| \dot{U}_{C\textsl{(ВО-1)}} - \frac{\dot{U}_{c\textsl{(ВО-2)}}}{\sqrt{3}} \right|
\end{equation}

\item
БНН по критерию пофазного сравнения напряжений вводится в работу программным ключом \textbf{«БНН-ф»}. 

\item
Неисправность в цепях напряжения определяется, если величина векторной невязки по одной из фаз превышает уставку \textbf{«Uср БНН-ф»}.

\item
В зависимости от схемы подключения дополнительной вторичной обмотки с помощью программных переключателей \textbf{Ua ВО-2»}, \textbf{«Ub ВО-2»} и \textbf{«Uc ВО-2»} можно задать, какой вывод «разомкнутого треугольника», соответствует фазе «звезды». Для предотвращения ложного срабатывания БНН-Ф при однофазном КЗ внутри контура заземления подстанции (в случае неправильного выполнения цепей заземления) предусмотрена блокировка по току нулевой последовательности, которая вводится программным ключом \textbf{«Контр 3I0 БНН-ф»}. Уставка блокировки по току нулевой последовательности не регулируется и равна 
Iном.  

\end{enumerate}


\item БНН по принципу КРБ-12М (БНН)
\begin{enumerate}


\item
В устройстве предусмотрен алгоритм БНН, повторяющий принцип построения КРБ-12М, который реагирует на напряжение небаланса $U_{\textsl{БНН}}$ , В, между основной и дополнительной вторичными обмотками.  

\item
Расчет $U_{\textsl{БНН}}$,В , зависит от схемы подключения цепей «разомкнутого треугольника». В таблице приведены уставки для БНН по сравнению напряжений «звезды» и «разомкнутого треугольника». 


%%%%%%%%%%%%%ТАБЛИЦА %%%%%%%%%%%%%%%%%%%
\small
\setlength{\extrarowheight}{0.05cm} 
\begin{longtable}
{|>{\centering\arraybackslash}m{6cm}|>{\centering\arraybackslash}m{4cm}|>{\centering\arraybackslash}m{2cm}|>{\centering\arraybackslash}m{4cm}|}
\caption{Перечень уставок алгоритма БНН по принципу КРБ-12М\hfill\vspace{-0.5\baselineskip}}\label{bnn:tbl1}\\ 
\hhline{|-|-|-|-|}
\rowcolor{gray!30}
Наименование уставки  & Диапазон   & Шаг & Значение по умолчанию  \\ 
\hhline{|-|-|-|-|}
\endfirsthead
\caption*{\hspace{3pt}\emph{Продолжение таблицы \ref{bnn:tbl}\hfill\vspace{-0.5\baselineskip}}} \\ % сделано по ГОСТ 2.105 п.6.8.7
\hhline{|-|-|-|-|}
\rowcolor{gray!30}
Наименование уставки  & Диапазон   & Шаг & Значение по умолчанию  \\ 
\hhline{|-|-|-|-|}
\endhead
\hhline{|-|-|-|-|}
\endfoot
\endlastfoot
\centering Схема подключения цепей разомкнутого треугольника & \centering Uни/Uик \\ Uнф/Uфк \\ Uни/Uиф/Uфк \\ Uнк & \centering - & \centering \arraybackslash Uни/Uиф/Uфк \\  
\hline
\centering Особая фаза ТН  & \centering А/В/С & \centering - & \centering \arraybackslash А \\  
\hline
\centering Направление векторов звезды и треугольника  & \centering Совпадает/не совпадает  & \centering - & \centering\arraybackslash Совпадает \\  
\hline
\end{longtable}
\setlength{\extrarowheight}{0.0cm} % возврат параметра высота ячеек в таблице
\normalsize    

\item
Схемы соединения обмоток, формулы для расчета $U_{\textsl{БНН}}$, векторные диаграммы приведены в приложении А. 

\item
Алгоритм, повторяющий принцип построения КРБ-12М, вводится в работу программным ключом \textbf{«ИО БНН»}. Порог срабатывания по небалансу напряжений задается уставкой \textbf{«Uср ИО БНН»}.

\end{enumerate}


\item Комбинированный БНН (Комб. БНН)
\begin{enumerate}

\item
Также предусмотрен алгоритм БНН, основанный на комбинированном принципе, контролирующий напряжения и токи. Алгоритм вводится программным ключом \textbf{«Комб БНН»}. Для выявления несимметричных повреждений в цепях напряжения (обрыв или короткое замыкание одной или двух фаз) контролируются следующие условия:

\begin{itemize}
\item   превышение напряжения обратной последовательности заданного значения (более $0,085\cdot U_{\textsl{ном}}$);

\item   отсутствие тока обратной последовательности. БНН блокируется при превышении отношения 
        тока обратной последовательности к току прямой последовательности заданной уставки \textbf{«I2/I1 комб 
        БНН»} – это необходимо для исключения излишнего срабатывания БНН при несимметричных коротких 
        замыканиях в силовой цепи. Рекомендуемым значением коэффициента является 0,2; 

\item   контроль отношения напряжения обратной последовательности к напряжению прямой 
        последовательности. БНН срабатывает только при превышении отношения заданной уставки 
        \textbf{«U2/U1 комб БНН»}. Рекомендуемым значением коэффициента является 0,2;

\item   отсутствие броска тока прямой последовательности. БНН блокируется при превышении приращения тока прямой последовательности заданной уставки \textbf{«dI1 комб. БНН»}.
\end{itemize}

\item 
Обрыв трех фаз как во вторичной цепи, так и на стороне ВН, определяется при одновременном выполнении следующих условий:

\begin{itemize}
    \item   скачкообразное снижение напряжения прямой последовательности более чем на 0,2$\cdot$Uном;
    \item   отсутствие броска тока прямой последовательности, превышающего заданную величину, задаваемую уставкой \textbf{«dI1 комб. БНН»}.
\end{itemize}

\item
При этом переход несимметричного повреждения в симметричное не приводит к возврату БНН, а также обеспечивается отсутствие ложного срабатывания при КЗ в силовой цепи.

\item
Сигнал обрыва трех фаз подхватывается энергонезависимым триггером при снижении напряжения прямой последовательности ниже уставки \textbf{«U1 мин БНН»}. Возврат цепи происходит после повторного повышения напряжения прямой последовательности более уставки \textbf{«U1 мин БНН»} (рекомендуемое значение уставки 0,1$\cdot$Uном). Если измерительные ТН установлены на линии, то дополнительно контролируется включенное положение выключателя.

\item
Для контроля одновременного исчезновения трех фазных напряжений используются три реле минимального напряжения в фазах А, В, С. Отсутствие напряжения определяется при снижении всех фазных напряжений меньше уставки \textbf{«Uф мин БНН»} в течение 5 секунд и наличии тока в защищаемой линии, определяемого по сработанному состоянию ИО УРОВ. Отсутствие напряжения действует на формирование сигнала обрыва трех фаз.

\end{enumerate}


\item Контроль обрыва нуля схемы «звезда»
\begin{enumerate}[label=\arabic{section}.\arabic{subsection}.\arabic{enumi}.\arabic*, labelsep=4pt, leftmargin=0em, itemindent=65pt, parsep=0pt]

\item
Логика БНН включает в себя контроль обрыва нуля схемы «звезда». Для выявления обрыва нулевого провода предусмотрено создание искусственной несимметрии путём подключения догрузочного резистора к входу фазы В в измерительных цепях напряжения. Обрыв нулевого провода определяется при превышении напряжения нулевой последовательности 0,085$\cdot$Uном и отсутствии появления тока нулевой последовательности, контролируемого уставкой \textbf{«3I0/I1 БНН»}. Данный алгоритм также реагирует на обрыв одной или двух фаз цепей напряжения. Ввод алгоритма осуществляется программным ключом \textbf{«Контр. Обрыв НП»}.

\item 
В логике БНН предусмотрен вход ФСУ \textit{«Фикс. Неиспр. ЦН»} для принудительного ввода состояния неисправности измерительных цепей напряжения.

\end{enumerate}


\item Контроль состояния автомата ТН 
\begin{enumerate}[label=\arabic{section}.\arabic{subsection}.\arabic{enumi}.\arabic*, labelsep=4pt, leftmargin=0em, itemindent=65pt, parsep=0pt]

\item
Для контроля состояния автомата ТН предусмотрен вход ФСУ \textit{«Автомат ТН»}. Сигнал может приниматься как от нормально открытых, так и от нормально закрытых блок-контактов. Выбор типа блок-контактов осуществляется программным переключателем \textbf{«Конт. Автомат ТН»}. Неудачная попытка восстановления напряжения при включении автомата ТН на КЗ не приводит к возврату БНН.

\item
Действие БНН осуществляется без выдержки времени на блокировку зависимых алгоритмов РЗА, а с выдержкой времени \textbf{«Тсигн неиспр. ЦН»} происходит действие на сигнализацию.

\end{enumerate}

\item
Оперативный ввод/вывод БНН производится с помощью БУ (блоков управления функциями) «Вывод БНН».  


\begin{figure}[H]
\centering
\includegraphics[width=1\textwidth,height=1\textheight,keepaspectratio]{\fbpath/261.1461 БНН/img/BNN.pdf}
\captionof{figure}{Функциональный блок БНН}\label{pict:BNN}
\end{figure}

\end{enumerate} 